\section{The protocol, its primitives, and CryptoVerif}\label{sec:background}

\subsection{HYBRID-AKE}\label{sec:protocol}

Two parties derive a 32-byte session key. Each authenticating party has a
long-term composite signature key pair $(\vk,\sk)$ and a name $\id$; a
verifier obtains its peer's $\vk$ out of band. The initiator generates a
fresh ephemeral KEM key pair $(\pke,\ske)$ for each session and erases
$\ske$ once the key is derived. The protocol runs in two modes:
\emph{responder-authenticated}, in two messages with an anonymous initiator,
and \emph{mutually authenticated}, in three. \Cref{fig:protocol} shows the
message flow.

Both parties compute the transcript hash and the session key as
\begin{align*}
  \THH &= \Hash\bigl(\LP(\SUITE, \id_I, \vk_I, \id_R, \vk_R, \pke, \ct)\bigr),\\
  K   &= \mathrm{HKDF}\bigl(\mathit{ikm} = \ssv,\ \mathit{salt} = \THH,\
         \mathit{info} = \mathsf{INFO} \,\|\, \texttt{0x20} \,\|\, \SUITE,\ L = 32\bigr),
\end{align*}
where $\Hash$ is SHA3-256~\cite{FIPS202}, HKDF~\cite{RFC5869,Kra10} is
instantiated with SHA-256, and $\LP$ concatenates its arguments, each
prefixed by its 4-byte big-endian length; \texttt{0x20} is a single space
byte. $\SUITE$ names the KEM and
signature profiles, and $\id_I,\vk_I$ are empty in responder-authenticated
mode. The responder signs $\THH$ under the context string
$\CTXR = \text{\texttt{"hybrid-ake/responder-auth/v1"}}$. In mutual mode the
initiator signs $\Hash(\LP(\THH, \sigma_R))$ under
$\CTXI = \text{\texttt{"hybrid-ake/initiator-auth/v1"}}$, which binds its
signature to the responder's. The initiator accepts $K$ only after verifying
$\sigma_R$; a mutual-mode responder accepts only after verifying $\sigma_I$.

\begin{figure}[t]
\centering
\small
\renewcommand{\arraystretch}{1.25}
\begin{tabular}{@{}l@{\hspace{4pt}}c@{\hspace{4pt}}l@{}}
\textbf{Initiator} $I$ & & \textbf{Responder} $R$\\[2pt]
\hline\\[-8pt]
$(\pke,\ske) \gets \mathsf{KeyGen}()$ & & \\
 & $\xrightarrow{\ m_1 \,=\, \LP(\SUITE,\ \id_I,\ \vk_I,\ \pke)\ }$ & \\
 & & check $\SUITE$\\
 & & $(\ct,\ssv) \gets \mathsf{Encaps}(\pke)$\\
 & & compute $\THH$\\
 & & $\sigma_R \gets \mathsf{Sign}(\sk_R, \THH, \CTXR)$\\
 & & $K \gets \KDF(\ssv; \THH)$\\
 & $\xleftarrow{\ m_2 \,=\, \LP(\id_R,\ \vk_R,\ \ct,\ \sigma_R)\ }$ & \\
$\ssv \gets \mathsf{Decaps}(\ske, \ct)$ & & \\
compute $\THH$ & & \\
require $\mathsf{Verify}(\vk_R, \sigma_R, \THH, \CTXR)$ & & \\
$K \gets \KDF(\ssv; \THH)$ & & \\
\multicolumn{3}{@{}l}{\emph{mutual mode only}, with $h = \Hash(\LP(\THH,\sigma_R))$:}\\
$\sigma_I \gets \mathsf{Sign}(\sk_I, h, \CTXI)$ & & \\
 & $\xrightarrow{\ m_3 \,=\, \LP(\sigma_I)\ }$ & \\
 & & require $\mathsf{Verify}(\vk_I, \sigma_I, h, \CTXI)$\\[2pt]
\hline
\end{tabular}
\caption{HYBRID-AKE. $\THH$ and $\KDF(\ssv;\THH)$ are as in
\cref{sec:protocol}; $\mathsf{KeyGen}$, $\mathsf{Encaps}$ and $\mathsf{Decaps}$
are the hybrid KEM's, and $\mathsf{Sign}$ and $\mathsf{Verify}$ the composite
signature's. In responder-authenticated mode, $\id_I$ and $\vk_I$ are empty
and $m_3$ is not sent.}
\label{fig:protocol}
\end{figure}

The protocol has one implementation, in Python on top of
pyca/cryptography~\cite{pyca}, and a symbolic analysis in Tamarin covering
the same goals, an ephemeral-key-reveal oracle and cross-suite downgrade
resistance (\cref{sec:related}).

\subsection{The hybrid KEMs}\label{sec:kems}

The KEM is one of two profiles.

\paragraph{X-Wing.} ML-KEM-768 with X25519~\cite{RFC7748}, as specified in
the Internet-Draft \ident{draft-connolly-cfrg-xwing-kem-11}~\cite{XWingDraft}. Public keys and
ciphertexts are concatenations, $\pk = \pk_M \| \pk_X$ and
$\ct = \ct_M \| \ct_X$, and the shared secret is
\[
  \ssv = \text{SHA3-256}\bigl(\ssv_M \,\|\, \ssv_X \,\|\, \ct_X \,\|\, \pk_X \,\|\, \mathsf{XWingLabel}\bigr),
\]
where $\ssv_M$ and $\ssv_X$ are the ML-KEM and X25519 shared secrets. The
combiner deliberately omits $\ct_M$ and $\pk_M$; the X-Wing paper justifies
this by proving that ML-KEM-768 is ciphertext second-preimage resistant
(C2PRI) when its hash functions are random oracles~\cite[Thm.~3]{XWing}. The
decapsulation key is a 32-byte seed, expanded with SHAKE-256 into the ML-KEM
seed and the X25519 scalar. (The published X-Wing paper puts the label first
in the hash input~\cite{XWing}; draft-11 puts it last, and we follow the
draft.)

\paragraph{ML-KEM-1024 + X25519.} X-Wing specifies no ML-KEM-1024 variant,
although the combiner it uses is also secure with
ML-KEM-1024~\cite{Starfighters}. For NIST Category~5 the implementation
instead uses a transcript-binding combiner in the style of~\cite{BBFGS19}:
\begin{align*}
  \ssv = \mathrm{HKDF\text{-}SHA256}\bigl(
    &\mathit{salt} = \text{SHA-256}(\text{\texttt{"hybrid-kem-v1"}}),\
     \mathit{ikm} = \ssv_M \,\|\, \ssv_X,\\
    &\mathit{info} = \LP(\ct_M, \ct_X, \pk_M, \pk_X)\bigr).
\end{align*}
It hashes both ciphertexts and both public keys. Its X25519 leg rejects an
all-zero Diffie--Hellman output in both directions. That follows RFC~9180's
rule for DHKEM recipients~\cite[\S7.1.4]{RFC9180}, although the profile is
not DHKEM.

\subsection{Composite signatures}\label{sec:composite}

\begin{sloppypar}
The signature is a composite of ML-DSA and EdDSA, following the IETF LAMPS
working-group draft \ident{draft-ietf-lamps-pq-composite-sigs-19}~\cite{Composite}.
We use its profiles \ident{id-MLDSA65-Ed25519-SHA512} and
\ident{id-MLDSA87-Ed448-SHAKE256}. To sign a message $M$ with context
$\mathit{ctx}$, both components sign the same message representative
\[
  M' = \mathsf{Prefix} \,\|\, \mathsf{Label} \,\|\, \mathrm{len}(\mathit{ctx}) \,\|\, \mathit{ctx} \,\|\, \PH(M),
\]
ML-DSA with its context set to $\mathsf{Label}$. $\PH$ is SHA-512 or
SHAKE256 with 64-byte output, depending on the profile. A composite
signature is the pair of component signatures, and it is valid only if both
verify.
\end{sloppypar}

\subsection{CryptoVerif}\label{sec:cv}

CryptoVerif~\cite{Bla08,CryptoVerif} proves security properties of protocols
and primitives in the computational model, by sequences of games. The
protocol and its environment are processes in a probabilistic calculus; we
use the ``oracles'' front-end, in which the adversary calls oracles that the
model defines. Assumptions on primitives are computational equivalences
$L \approx_p R$ between two sets of oracles, collected in libraries of macros
such as \ident{IND_CCA2_KEM}, \ident{UF_CMA_proba_signature},
\texttt{PRF}, \ident{ROM_hash_large} and \texttt{StDH}. A proof applies
these equivalences (\texttt{crypto} steps) and syntactic transformations
(\texttt{insert}, \texttt{replace}, \texttt{SArename}, \texttt{simplify},
\ldots). Each transformation is checked by the tool; in particular, a
\texttt{replace} succeeds only if CryptoVerif can prove the two terms equal.
The proof ends when the property can be read off the final game. CryptoVerif
then reports the adversary's advantage as a function of the assumptions'
probabilities, of the running times, and of the number of calls to each
oracle.

CryptoVerif proves secrecy in the real-or-random sense, with several test
queries; correspondences, including injective ones; and indistinguishability
of two processes. We use all three. CryptoVerif's post-quantum
library~\cite{BJ24} contains only assumptions that have instantiations secure
against quantum adversaries; it has no random oracles and no
Diffie--Hellman assumptions. A proof that uses only this library is
post-quantum sound.
